\scanpage{128}

Ce serait le moment d'ailleurs de noter que l'on a

\[
Z'\subset G \tag*{(58)}
\]

car on voit \deleted{que} pour $(s,t)\in S^2$

\[
\{s,t\}\in A_{\deleted{\unsure{i}}i_0}\Longleftrightarrow
\{s,\lambda t\}\in A_{(\lambda i_0)}\Longleftrightarrow
\{\lambda s,\lambda t\}\in A_{\underbrace{\lambda.(\lambda i_0)}_{\lambda^2 i_0}}
\]

d'où.

\[
\underbrace{\lambda(A_{i_0})}_{\substack{\text{transport}\\\text{de structure}}}
=A_{\underbrace{\lambda^2.(i_0)}_{\substack{\text{opération}\\\text{de }Z'\text{ sur}\\\Pi\text{ définie par (45)}\\\text{p. }61'}}}, \tag*{(59)}
\]

% slanted
\margin{attention, tournant dangereux et \deleted{\unsure{important}}\inserted{\unsure{mystérieux}} car $\operatorname{det}.(\lambda\,\mathrm{id}_M)=\lambda^2$ (\unsure{$\#$|$\ne$} $\lambda$ en \deleted{\unsure{tant que|général}} que fonction en $\lambda$!).}

et l'opération de $G$ sur $Z'$ est induite par autom. intérieurs. Ceci dit, on va poser

\[
\begin{aligned}
(T_0\,7)\quad&Z'\ \deleted{\text{\unsure{centralise}}}\inserted{\text{est central}}\text{ dans }G\quad\text{(NB. il contient}\\
&\text{de }\deleted{\text{\unsure{toutes}}}\text{ façons le centre de }G,\text{ donc }T_0\,7\\
&\text{signifie aussi}\\
&\hspace{2em}Z'=\operatorname{Cent}(G)\unsure{,|)}
\end{aligned}\tag*{(60)}
\]

qui équivaut aussi au fait que $G$ opère sur $\Pi$ par translations, i.e. (56) trivial i.e. (57) donné par un hom de groupes

\[
\deleted{\unsure{Z}}G\xrightarrow{\delta}Z', \tag*{(61)}
\]

de sorte qu'on a finalement une suite exacte. canonique
% archive sheet number: 127
